\documentclass[a4paper]{article}
\usepackage[pdftex]{graphicx}
\usepackage[utf8]{inputenc}
\usepackage{enumerate}
\usepackage{siunitx}
\sisetup{locale=DE} 
\usepackage{href-ul}
\hypersetup{
	colorlinks=true,
	linkcolor=blue,
	urlcolor=blue}
\usepackage{icomma}
\usepackage{amssymb}
\usepackage{geometry}
\geometry{a4paper, top=15mm, left=15mm, right=15mm, bottom=15mm,
	headsep=10mm, footskip=12mm}
%Source: img_0185

% Start the document
\begin{document}
%Titel
{\bf Aufgaben Leiterkennlinien }

\begin{enumerate}[1.]

%Aufgabe 1
\begin{minipage}{0.65\textwidth}
{\bf \item In einem Experiment wird die Stromstärke $I$ in Abhängigkeit von der Spannung $U$ }für einen Konstantandraht, eine Bleistiftmine aus Graphit und einen Eisendraht untersucht. Es ergibt sich das nebenstehende Diagramm. 
\end{minipage}
\hfill
\begin{minipage}{0.3\textwidth}
\includegraphics[width=4 cm]{leiken185}
\end{minipage}

\begin{enumerate}[a)]
\item  Was kann man über die Widerstandswerte der drei Leiter bei steigender Spannung aussagen? Ordnen Sie den drei Graphen die entsprechenden Leiter zu. Wie bezeichnet man solche Graphen allgemein? 
\item Welches der drei Materialien leitet bei einer Spannung von $\SI{4,0}{\volt}$ am besten?
Berechnen Sie für diesen Leiter den dazugehörigen Widerstandswert.
\item Erklären Sie mit Hilfe des Teilchenmodells den Verlauf der Kennlinie von Leiter 2.
\item Erläutern Sie Begriffe  PTC und NTC und was ein Kalt- und was ein Heißleiter ist.
\end{enumerate}

%Aufgabe 2
{\bf \item In einem Versuch wird die Stromstärke $I$ in Abhängigkeit von der Spannung $U$ für zwei verschiedene elektrische Leiter untersucht.} Dabei ergeben sich folgende Messwerte:

\begin{tabular}{ccccccccc}
& $U$ in $\SI{}{\volt}$ & 0 & 1,8 & 3,6 & 5,4 & 7,2 & 9,0 & 10,8 \\
\hline
Leiter 1 & $I$ in $\SI{}{\ampere}$ & 0 & 0,31 & 0,55 & 0,74 & 0,84 & 0,93 & 0,96 \\
Leiter 2 &  $I$ in $\SI{}{\ampere}$ & 0 & 0,18 & 0,37 & 0,53 & 0,71 & 0,92 & 1,06
\end{tabular}

\begin{enumerate}[a)]
\item Werten Sie die Messtabelle für Leiter 1 numerisch aus und formulieren Sie das Versuchsergebnis.
\item Stellen Sie für Leiter 2 die Stromstärke  in Abhängigkeit von der Spannung in einem Diagramm graphisch dar, formulieren und begründen Sie das Versuchsergebnis.
\item Welche Aussagen kann man mit Hilfe der numerischen und der gra\-phi\-schen Auswertung über die Widerstandswerte der beiden Leiter treffen?
\end{enumerate} 

%Aufgabe 3
\item {\bf Elektrischer Widerstand }

\begin{enumerate}[a)]
\item Der elektrische Widerstand einer Glühlampe soll experimentell bestimmt werden. Dazu wird eine anliegende Spannung von $\SI{6,0}{\volt}$ sowie eine Stromstärke von $\SI{5,0}{\ampere}$ gemessen. Berechne den Widerstand des Lei\-ters.
\item Der Widerstand eines Leiters beträgt $\SI{50}{\ohm}$. Die Stromstärke im Draht wird zu $\SI{380}{\milli\ampere}$ bestimmt. Welche Spannung liegt an?
\item Eine Glühlampe mit der Leistung $P=\SI{40}{\watt}$ ist an das Haushaltsnetz ($U=\SI{230}{\volt}$) angeschlossen. Berechne die Betriebsstromstärke und den Betriebswiderstand.
\item Eine Stromstärke von $\SI{15}{\milli\ampere}$ kann für den Menschen bereits tödlich sein. Unter ungünstigen Bedingungen beträgt der elektrische Widerstand des menschlichen Körpers etwa 
$\SI{1,0}{\kilo\ohm}$. Welche Spannung kann unter diesen Umständen bereits lebensgefährlich sein?
\end{enumerate}
\end{enumerate} 

\fbox{
	\begin{minipage}{0.5\textwidth}
		Zur Lösung bitte \href{https://www.okuyakl.de/phys/p8widL185/ll185.pdf}{hier klicken} oder den QR-Code scannen.\\
	Weitere Arbeitsblätter gibt es unter 
	
	\href{https://www.okuyakl.de}{www.okuyakl.de}
	\end{minipage}
	\hfill
	\begin{minipage}{0.4\textwidth}
		\includegraphics[width=1.5 cm]{../../viecher/zwe03}
		\includegraphics[width=3 cm]{qr185}
		\includegraphics[width=2 cm]{../../viecher/afanticon1}
		
	\end{minipage}}

\end{document}%L Ö S U N G -----------------------------------------
